\documentclass[10pt,a4paper]{amsart}
\usepackage[margin=19mm]{geometry}
\usepackage{amsmath,amssymb,mathtools}
\usepackage[scr=rsfs,cal=euler]{mathalfa}
\usepackage{xfrac}
\usepackage[unicode,hidelinks,bookmarks=false]{hyperref}
\newcommand{\ie}{i.e.\ }
\newcommand{\cf}{cf.\ }
\newcommand{\resp}{respectively\ }
\newcommand{\Subsection}{\subsection}
\providecommand{\nobreakdash}{\nobreak\hbox{-}}
\setlength{\emergencystretch}{2em}
\multlinegap0pt
\usepackage{bm}
\usepackage{bbm}
\usepackage{dsfont}
\usepackage{xspace}
\usepackage{cancel}
\usepackage{mathtools}
\usepackage{amsthm}
\makeatletter
\@ifundefined{theorem}{\newtheorem{theorem}{Theorem}}{}
\@ifundefined{proposition}{\newtheorem{proposition}[theorem]{Proposition}}{}
\makeatother
\newcommand{\C}{\mathbb{C}}
\newcommand{\Ox}{\mathcal{O}}
\newcommand{\R}{\mathcal{R}}
\newcommand{\Z}{\mathbb{Z}}
\newcommand{\p}{\mathbb{P}}
\title[On the Cox rings of some hypersurfaces]{On the Cox rings of some hypersurfaces}
\author{Andrew Pollock, Atsushi Ito, Balazs Szendroi}
\date{Source: arXiv:2504.07591v2 (CC BY 4.0)}
\begin{document}
\maketitle

\noindent\emph{Source and licence.} On the Cox rings of some hypersurfaces. Version arXiv:2504.07591v2 (CC BY 4.0), distributed under the Creative Commons Attribution 4.0 International licence, as recorded in the arXiv licence metadata for this exact version and shown on the abstract page https://arxiv.org/abs/2504.07591v2. Full rights notice: licenses/arxiv-2504.07591-CC-BY-4.0.txt.\par\smallskip

\noindent\textit{Editorial scope.} Counted unit: the Proposition whose source label is \texttt{prop:linear} in the article (source0.tex lines 802--804, source label \texttt{prop:linear}), delivered together with the article's own complete original proof of it (source0.tex lines 806--857, one \texttt{proof} environment of 52 lines). No mathematics is written, repaired or re-derived: statement and proof are the article's own text line for line. Required context retained uncounted: prop:z\_k (lines 425--433); eq:linear (lines 775--777). Every definition, display and label of those lines is retained unchanged. Omitted from this package and not redistributed: 1--424, 434--774, 778--801, 805--805, 858--1114. Nothing inside a retained excerpt is omitted or altered: every retained body is delivered exactly as the article's lines are written, so no omission receipt of any kind accompanies this package. Citations: the retained excerpts carry no citation command at all, so no bibliography entry of the article is transcribed and no reference is redistributed. Reference adaptation: none was needed and none was made; every reference inside the delivered unit resolves inside it, so no unresolved reference or citation is produced. Printed slips kept literal: no printed word, symbol, hypothesis or number of the counted statement or of its proof was altered. Layout and class: the article's own class is \texttt{amsart} with packages including geometry, inputenc, babel, amssymb, amsmath, amsthm, amsfonts, array, comment, enumitem, hyperref, mathrsfs, tikz, tikz-cd; the wrapper is the lane's pinned \texttt{amsart} editorial wrapper at 10pt on A4, which restates the theorem-like environments the retained text needs and adds the article's own macro definitions that the retained text needs (\texttt{\textbackslash C}, \texttt{\textbackslash Ox}, \texttt{\textbackslash R}, \texttt{\textbackslash Z}, \texttt{\textbackslash p}), with extra wrapper packages (bm, bbm, dsfont, xspace, cancel, mathtools). The delivered numbering is the unit's own. Assets: no retained span contains a figure environment, a graphics-inclusion command, a PDF-page inclusion, a table or a TikZ picture; no image file is copied into this package, the package declares no asset dependency and no image file is redistributed with it. Licence: version arXiv:2504.07591v2 of this article is distributed under the Creative Commons Attribution 4.0 International licence, as recorded in the arXiv licence metadata for this exact version and shown on the abstract page https://arxiv.org/abs/2504.07591v2. A full rights notice is delivered at licenses/arxiv-2504.07591-CC-BY-4.0.txt. Characters: every retained excerpt is pure ASCII, so the delivered engine, XeLaTeX with its default Latin Modern text font, reports no missing character.
\begin{proposition}\label{prop:z_k}
\begin{enumerate} \item[(i)] There exists sections $z_1, \dots, z_d \in \R(Z)_{(-1,e)}$, linearly independent over $R$, that satisfy the $(d+1)$ equations
\begin{equation}
    x_1 z_1 + g_0 = x_1 z_{2} + g_{1} - x_0 z_{1} = \dots = x_1 z_{d} + g_{d-1} - x_0 z_{d-1} = g_d - x_0 z_d = 0.
    \label{eq:coxz}
\end{equation}
\item[(ii)]The subalgebra $\R(Z)_{z_k}$ of $\R(Z)$ generated by $x_0, x_1, y_0, \dots, y_n, z_1, \dots, z_d$ contains $\R(Z)_{(-a,b)}$ for all $-a \geqslant -1, b \in \Z$.\end{enumerate}
\label{prop:coxz1}
\end{proposition}

\begin{equation}\label{eq:linear}
       x_0 w = \lambda_0 z_1 + z_2 +\lambda_1 z_3 +\lambda_2 z_4, \quad x_1 w = z_1 + \lambda_1 z_2+ \lambda_2 z_3 +\lambda_3 z_4  .
\end{equation}

\begin{proposition}\label{prop:linear}
         For general constants $\lambda_j\in\C$, the $\Z^2$-graded Cox ring $\R(Z)$ of the hypersurface $Z \subset \p^1 \times \p^3$ is isomorphic to the polynomial ring $\C[x_0,x_1,w,z_3,z_4]$.    
\end{proposition}

\begin{proof}
Let $r \in H^0(\Ox_{\p^1 \times \p^3}(4,1)) =H^0(\Ox_{\p^1}(4)) \otimes H^0( \Ox_{\p^3}(1)) $ be the section which defines $Z$.
Then $r$ induces an exact sequence
\begin{align*}
0 \to \Ox_{\p^1}(-4) \xrightarrow{r} H^0( \Ox_{\p^3}(1)) \otimes \Ox_{\p^1} \to \mathcal{E} \to 0
\end{align*}
on $\p^1$
and we have $Z =\p_{\p^1}(\mathcal{E})\subset \p^1 \times \p^3$.
We note that $\mathcal{E}$ is locally free of rank $3$
since $r$ is general.
Let $\mathcal{E} \simeq \Ox_{\p^1} (a) \oplus \Ox_{\p^1} (b) \oplus \Ox_{\p^1} (c)$ with $a+b+c=4$.
Since $\mathcal{E}$ is a quotient of a trivial bundle, we have $a,b,c \geq 0$.
If $a =0$, then $r$ must be contained in $H^0(\Ox_{\p^1}(4)) \otimes V'$ for
$V' =\ker H^0(\Ox_{\p^3}(1)) \to H^0(\Ox_{\p^1}(a))=\C$,
which contradicts the generality of $r$.
Hence $a,b,c \geq 1$ and $\mathcal{E} \simeq  \Ox_{\p^1} (1) \oplus \Ox_{\p^1} (1) \oplus \Ox_{\p^1} (2)$. 


Since the embedding $Z =\p_{\p^1}(\mathcal{E})\subset \p^1 \times \p^3$ over $\p^1$ is induced by 
$H^0( \Ox_{\p^3}(1)) \otimes \Ox_{\p^1} \twoheadrightarrow \mathcal{E} $,
the tautological line bundle $\Ox_{\mathcal{E}}(1)$ on $ \p_{\p^1}(\mathcal{E})$ is identified with $\Ox_Z(0,1)$,
where $\Ox_Z(a,b) \coloneqq \Ox_{\p^1\times \p^3}(a,b)|_Z $.
Since $\mathcal{E}(-1) \simeq 
\Ox_{\p^1} \oplus \Ox_{\p^1} \oplus \Ox_{\p^1} (1)$, 
the morphism 
%\balazs{Balazs: I would write this as
%\[\mu : Z=\p_{\p^1}(\mathcal{E}(-1)) \to P = \p H^0(\Ox_Z(-1,1)) \cong\p^3\]
%}
\[\mu : Z=\p_{\p^1}(\mathcal{E}(-1)) \to P = \p H^0(\Ox_Z(-1,1)) \cong\p^3\]
induced by $|\Ox_Z(-1,1)|$ is the blow-up along a line $l$.
We note that $P$ is different from the second factor of $\p^1 \times \p^3$.


Let 
\begin{align*}
D_i =(x_i=0), \quad D'_k=(z_k =0), \quad E=(w=0) 
\end{align*}
be divisors on $Z$.
Since $w$ is a non-zero section of $H^0( \Ox_Z(-2,1)) = H^0(\mathcal{E}(-2)) \simeq \C$,
$E$ is the exceptional divisor of the blow-up $\mu$.

Since $z_1,z_2,z_3,z_4$ are linearly independent over $R$ by Proposition \ref{prop:z_k},
so are over $\C$.
Hence $z_1,z_2,z_3,z_4$ form a basis of $ H^0(\Ox_Z(-1,1))  =H^0(\mathcal{E}(-1)) \simeq \C^4$.
By (\ref{eq:linear}), $x_0w,x_1w,z_3,z_4$ also form a basis of $ H^0(\Ox_Z(-1,1)) =H^0(P, \Ox(1))$. We regard $P=\p^3$ as a toric variety by this basis.
Let $H_i=(x_iw=0) \subset P$ be the torus invariant planes for $i=0,1$.
Since $\mu^*H_i=D_i+E$, the line $l =\mu(E)$ is $H_0 \cap H_1$, which is torus invariant.
Hence $\mu$ is a toric morphism and $Z$ is a toric variety as well.
Then torus invariant prime divisors of $Z$ are $D_0,D_1,E, D'_3,D'_4$.
Since the Cox ring of a smooth projective toric variety is a polynomial ring whose variables correspond to torus invariants prime divisors,
we conclude that $\R(Z) \simeq \C[x_0,x_1,w,z_3,z_4]$.
\end{proof}

\end{document}
